\documentclass[../Main.tex]{subfiles}
\begin{document}
\section{Conclusion}

This thesis presented the design and implementation of FloodHelper, a disaster assistance platform that integrates real-time location sharing, emergency rescue coordination, and charity campaign management. The system was developed to improve communication and coordination among citizens, rescuers, benefactors, authorities, and administrators during disaster situations through a unified mobile and web-based platform.

Based on the requirements analysis, the system was designed using a hybrid software architecture that combines the Client–Server model with the Publish–Subscribe communication pattern. The mobile application provides functionalities for user management, friend management, location sharing, emergency distress broadcasting, rescue coordination, and charity campaign participation, while the web application supports administrative and authority-related operations. The backend system, together with the MQTT communication service, enables both conventional request-response processing and real-time data synchronization.

The implementation successfully realizes the major functional requirements identified in the analysis phase. Users can register and manage their accounts, establish friendships, share their real-time locations with authorized friends, broadcast distress signals, participate in rescue activities, and contribute to charity campaigns. Authorities are able to review role applications and campaign requests, while administrators manage user accounts and maintain overall system operation. The integration of these modules provides a unified platform that supports multiple disaster-related activities within a single application.

Overall, the implemented system demonstrates the feasibility of integrating social networking, emergency response, and charity campaign management into a unified disaster assistance platform. The adopted system architecture provides a modular and maintainable foundation that can be further extended to support additional functionalities and larger deployment environments in future development.

% \subsection{Comparison with Similar Systems}
% In this graduation thesis, the developed system, FloodHelper, has been compared against existing communication and tracking systems widely used in emergency or location-sharing contexts, specifically instant messaging applications (e.g., Zalo), web-based SOS platforms, and social location applications (e.g., Whoo). Compared to these platforms, the proposed system demonstrates significant architectural and functional advancements. Unlike standard instant messaging systems or web-based SOS systems which require manual activation within a session or fail entirely when the network connection drops, FloodHelper leverages the lightweight MQTT protocol and background isolates to capture and cache the user's last known location. This allows rescuers to establish a search perimeter based on the victim's last known coordinates even when the device is completely offline. Furthermore, FloodHelper facilitates direct, ward-based rescue coordination rather than relying on manual contact sharing or browser-based platforms that lack rescuer collaboration tools.

% When compared to social location-tracking applications such as Whoo, which are designed exclusively for sharing coordinates among friends, the primary strength of FloodHelper lies in its ability to directly connect victims in distress with professional and volunteer rescue teams. While Whoo successfully implements background location tracking and last-location caching, it is fundamentally a social networking platform that lacks any mechanism to transmit SOS signals or alert emergency services. FloodHelper, on the other hand, structures location data on a ward-level basis to link victims directly with local rescue units, enabling real-time dispatching and coordination on an official administrative level. Furthermore, while social applications offer no support for post-disaster recovery, FloodHelper integrates a transparent charity campaign lifecycle module to manage donations and verify physical relief distribution, providing a unified solution that addresses both immediate rescue and long-term relief management.

% \subsection{Summary of Achievements}
% Throughout the development of this graduation thesis, designed and implemented functionalities have been successfully consolidated to serve two primary, interconnected goals: emergency rescue coordination and transparent charity management. Regarding the first goal of resilient emergency rescue and location tracking, we have successfully designed and built a mobile location-sharing and distress-signalling subsystem. This Flutter-based application allows victims to broadcast SOS alerts containing vital context, such as vulnerable demographics and resource needs, which are automatically prioritized and sorted for rescuers operating at the ward level. To support this under unstable network conditions, the communications pipeline is optimized using a Python-based MQTT Worker to handle high-frequency coordinate streams, achieving reliable offline last-location storing.

% For the second goal of transparent charity and donation management, the system digitizes the traditional relief lifecycle into a strictly monitored, state-driven workflow. Benefactors can create charity campaigns that undergo a strict vetting process by local Authorities via a dedicated NestJS backend and Web Authority Desk. Once approved, the system automates state transitions from donating to distributing, leveraging the VietQR API to track financial contributions in real-time. To ensure absolute public accountability, the system visually maps the benefactor's real-time geographical coordinates to verify the physical distribution of goods, and mandates the upload of official bank statements for final community auditing, restoring trust in post-disaster charity activities.

% \subsection{Key Contributions}

% In addition to delivering a comprehensive system that supports both emergency rescue operations and charitable campaign management, this project provides two notable technical contributions that directly address critical challenges encountered in real-world deployment scenarios. The first contribution is the design and implementation of an MQTT topic architecture, together with optimized publish/subscribe communication rules specifically tailored for location sharing and rescue coordination in environments with unstable Internet connectivity. The second contribution is the modeling and management of the lifecycle states of charitable campaigns, establishing a digital governance framework that helps prevent fraudulent activities. The following sections provide a detailed analysis of these contributions.

% \subsubsection{MQTT Topic Architecture and Communication Mechanism for Emergency Rescue Operations}

% The first major technical contribution lies in the messaging architecture, which was designed to overcome the limitations of weak Internet connectivity commonly found in flood-affected areas. The diagram in the figure \ref{fig:Fig613a} illustrates the topic network developed for location-sharing functionality among connected users (e.g., friends and family members). Rather than establishing direct communication between end-user mobile devices (Client A communicating individually with Clients B, C, and D), which may result in increased packet loss and inefficient bandwidth utilization, the system employs an MQTT Client operating as an MQTT Worker within a strong Internet connection environment to serve as an intermediary coordinator. When Client A publishes its current coordinates (\textit{myCurrentLocation}) to the \texttt{current-location} topic, the MQTT Worker receives, processes, and redistributes the location data by publishing it to personalized topics for each recipient, such as \texttt{A/to\_B/last-Location} and \texttt{A/to\_C/last-Location}. A key innovation of this design is the utilization of MQTT Retained Messages for destination topics. As a result, even if Clients B, C, and D temporarily lose network connectivity, they immediately receive the latest available location of Client A upon reconnection, without waiting for the next update cycle. This mechanism ensures the continuity and reliability of location information dissemination.

% \begin{figure}[H]
%     \centering
%     \includegraphics[width=\textwidth]{Figure/TopicNetwork/TopicForFriend.png}
%     \caption{Topic network for location sharing among friends}
%     \label{fig:Fig613a}
% \end{figure}

% Building upon the basic location-sharing mechanism, the topic architecture was further extended to support bidirectional coordination between Victims and Rescue Teams, as illustrated in the second diagram. When a victim (Client A) issues a distress signal (\textit{signalCommand}) through the \texttt{signal} topic, the signal is not indiscriminately broadcast across the network. Instead, it is routed directly to the MQTT Worker. The Worker synchronizes the information with the central backend system, records the incident in the database, and subsequently publishes a notification command (\textit{notificationCommand}) through the \texttt{rescuer-common} topic to the devices of rescue personnel (Rescuers B and C). Upon receiving the dispatch notification, rescue devices begin periodically publishing their current coordinates (\textit{rescuerLocation}) to the \texttt{rescuer-location} topic. On the victim side, the device only needs to subscribe to this topic to continuously receive updated rescue team locations, providing valuable situational awareness and psychological reassurance during emergency situations.

% \begin{figure}[H]
%     \centering
%     \includegraphics[width=\textwidth]{Figure/TopicNetwork/TopicForVictimAndRescuer_OldVersion.png}
%     \caption{Topic network for location sharing between victims and rescuers}
%     \label{fig:Fig613b}
% \end{figure}

% Although the proposed architecture effectively establishes a closed-loop communication workflow, the current implementation (as depicted in the second diagram) still exhibits a significant limitation regarding geographical information filtering. Under the existing topic-routing model, rescue personnel receive distress notifications from all victims nationwide, regardless of their geographical proximity. This limitation is not only impractical from an operational coordination perspective but also leads to inefficient consumption of network bandwidth, which is often a scarce resource in disaster-stricken areas. To address this issue, a technical enhancement has been identified that involves restructuring the topic hierarchy through the incorporation of geographical identifiers (e.g., \texttt{ward\_id}) as topic suffixes. This approach would isolate and partition communication flows according to specific administrative regions, thereby enabling location-aware notification routing. The proposed solution is discussed in greater detail in the \textit{Future Work} section of this project.

% \subsubsection{Charity Campaign Lifecycle State Management}

% While the MQTT-based communication architecture focuses on ensuring reliable rescue coordination under unstable network conditions, the second major technical contribution addresses a different challenge: establishing transparency and accountability in charitable activities. To achieve this objective, the system models each charity campaign as a finite-state machine (FSM), as illustrated in the figure \ref{fig:Fig613c}. By defining explicit states and controlled transitions, the system ensures that every campaign follows a predefined lifecycle, thereby reducing opportunities for fraud and improving operational governance.

% \begin{figure}[H]
%     \centering
%     \includegraphics[width=\textwidth]{Figure/StateMachine/CharityCampaignState.png}
%     \caption{Finite state machine diagram for charity campaign}
%     \label{fig:Fig613c}
% \end{figure}

% The lifecycle begins with the \textit{Created} state, which is entered immediately after a beneficiary creates a new campaign. At this stage, campaign information is still editable. Once the beneficiary submits the campaign for review, the campaign transitions to the \textit{Pending} state, indicating that it is awaiting verification. If the beneficiary wishes to modify any information during the review process, the campaign can return to the \textit{Created} state for further editing and resubmission.

% The transition from \textit{Pending} is governed by an automated verification process. When the system determines that all required information is valid and complete, the campaign is promoted to the \textit{Approved} state. Conversely, if the submitted information is found to be invalid, incomplete, or inconsistent, the campaign transitions to the terminal \textit{Rejected} state.

% Once approved, the campaign progresses automatically according to predefined scheduling milestones. Upon reaching the donation start date, the campaign enters the \textit{Donating} state, during which donors are allowed to contribute funds. When the distribution date is reached, the campaign transitions to the \textit{Distributing} state, indicating that collected resources are being allocated to beneficiaries. Finally, when the campaign reaches its designated end date, it moves to the \textit{Finished} state, marking the successful completion of all fundraising and distribution activities.

% In addition to the normal workflow, the state model incorporates governance mechanisms to ensure campaign transparency. During the \textit{Donating} and \textit{Distributing} phases, the responsible authority is granted the ability to review campaign activities and supporting information. If irregularities, inconsistencies, or suspicious behaviors are identified during the review process, the authority may transition the campaign to the \textit{Suspended} state. This action temporarily halts all fundraising and distribution activities until the identified issues have been investigated and resolved.

% A notable advantage of this state-machine-based approach is that it transforms campaign management from a collection of loosely controlled operations into a structured and auditable process. Every state transition is triggered by a clearly defined business event, such as beneficiary submission, automated validation, scheduled milestones, or fraud detection. Consequently, the system establishes a digital governance framework that enhances transparency, minimizes human error, and significantly reduces the risk of fraudulent charitable campaigns.

% \subsection{Limitations of the Work}
% Despite the successful implementation and validation of the primary rescue and charity features in the FloodHelper ecosystem, several limitations and uncompleted modules remain to be addressed:
% \begin{itemize}
%     \item Lack of Dynamic Weather Overlay: Although the map interface supports active pins for victims, rescuers, and campaign locations, it does not currently integrate dynamic meteorological layers. Users cannot view live weather maps (e.g., displaying temperature levels, rainfall density, or wind patterns), which restricts their ability to anticipate flood progressions.
%     \item Unlocalized Coordinate Sharing: While SOS signals are routed specifically to rescuers in the victim's ward, the standard real-time location-sharing module does not support ward-based partitioning. Non-emergency coordinates are shared broadly across friends rather than selectively with the local rescue team in the user's specific ward or commune.
%     \item Missing System Administrator Dashboard: While the web administration platform already includes basic user management functions for the Admin role, it lacks a comprehensive dashboard for tracking statistics and system telemetry. Global system administrators cannot easily view aggregated statistics on active charity campaigns, donation flows, rescuer performance, or rescue operations.
%     \item Lack of Integrated Real-Time Communication: Although the PostgreSQL database schema is fully designed to support chat rooms, messages, and room members, the current frontend mobile and web clients do not yet implement the user interface and service endpoints for real-time messaging between friends and rescue teams.
% \end{itemize}

% \subsection{Lessons Learned}
% The process of planning, designing, and implementing the full-stack FloodHelper platform has yielded several critical lessons in software development and architectural design:
% \begin{itemize}
%     \item Resiliency under Harsh Network Environments: Building real-time location streaming applications for emergency contexts underscored the importance of protocol selection. Utilizing lightweight MQTT topic rules and message retention proved far more reliable and efficient than heavy HTTP polling loops under low-bandwidth networks.
%     \item State Machines for Strict Business Rules: Modeling the charity campaign lifecycle as a Finite State Machine (FSM) demonstrated that representing complex systems through formal states prevents illegal state transitions (such as distributing funds before a campaign is approved). It also drastically simplified database modeling and audit trail verification.
%     \item Handling Mobile OS Platform Constraints: Encountering background task terminations by Android's power-saving configurations illustrated that production-ready software requires deep optimization at the OS-level. Implementing native foreground service wrappers is essential for keeping safety-critical tracking active.
%     \item Importance of Decoupled Multi-Platform Architecture: Developing four separate components (Flutter app, NestJS backend, Python MQTT worker, and PostgreSQL database) highlighted the necessity of rigid interface contracts. Standardizing DTO validation rules early in the design phase allowed parallel development and facilitated seamless end-to-end integration testing.
% \end{itemize}

\section{Future work}

Several directions can be considered to further improve the FloodHelper system.

First, future development may extend the mobile application to additional platforms, such as iOS, allowing a wider range of users to access the system.

Second, the real-time communication mechanism can be enhanced to better support unstable network conditions by incorporating offline synchronization or delay-tolerant communication techniques for disaster scenarios.

Third, the charity campaign module can be integrated with online payment gateways to simplify the donation process and provide a more convenient user experience.

Finally, future work may include comprehensive performance, scalability, and usability evaluations involving larger numbers of users and more realistic disaster simulation scenarios. Such evaluations would provide additional evidence of the system's effectiveness and support its deployment in practical environments.

% Based on the achievements obtained, the proposed system has successfully established the fundamental closed-loop workflows required to support its two primary objectives: emergency rescue coordination and charitable donation management. Nevertheless, in order to effectively handle large-scale data workloads, operate reliably under severely disrupted infrastructure conditions, and provide a more comprehensive ecosystem for real-world deployment, several future development directions are proposed. The following section outlines a roadmap for addressing the limitations of the current communication architecture while introducing new enhancements aimed at improving the system from multiple perspectives. These future developments focus on optimizing the distribution of geospatial data, strengthening offline information persistence through Delay-Tolerant Networking (DTN) technologies, supporting risk assessment and disaster preparedness through Weather Map integration, and completing the automated financial reconciliation mechanism based on VietQR.

% \subsection{Restructuring the MQTT Topic Network Based on Geospatial Partitioning}

% The highest-priority task for improving the rescue coordination functionality is addressing the issue of excessive message dissemination and bandwidth waste discussed in the previous section. To overcome this limitation, the MQTT routing architecture is proposed to be redesigned based on a geospatial partitioning mechanism, as illustrated in the figure \ref{fig:Fig621a}.

% \begin{figure}[H]
%     \centering
%     \includegraphics[width=\textwidth]{Figure/TopicNetwork/TopicForVictimAndRescuer_NewVersion.png}
%     \caption{Topic network after improvement}
%     \label{fig:Fig621a}
% \end{figure}

% Specifically, the generic rescue coordination topics used in the current implementation will be refined by appending administrative area identifiers at the ward/commune level (\texttt{<ward\_id>}). Communication flows will therefore be partitioned into region-specific topics, such as \texttt{rescuer/common/<ward\_id>}, \texttt{rescuer/handle/<ward\_id>}, and \texttt{rescuer-location/<ward\_id>}. To enable this mechanism to operate automatically, a dedicated Geospatial Service will be developed and integrated directly into the MQTT Client component acting as the MQTT Worker. Whenever raw GPS coordinates are received, the service will perform reverse geocoding to determine and extract the corresponding \texttt{ward\_id} associated with the location.

% In addition, to ensure that a victim's device (e.g., Client A) remains synchronized with the rescue network of its current region, a new context-awareness topic named \texttt{current-ward/A} will be introduced. Through this topic, Client A will always be informed of the administrative partition in which it is currently located. Based on this information, the application can automatically subscribe to the appropriate rescue coordination topics containing the corresponding \texttt{<ward\_id>} suffix. The completion of this architecture will establish a highly effective geospatial filtering mechanism. Rescue teams operating within a particular region will only receive distress signals and coordination requests originating from the same administrative area. This approach not only minimizes unnecessary network traffic and conserves critical bandwidth resources but also significantly improves the practicality, scalability, and operational efficiency of the rescue coordination system in real-world disaster scenarios.

% \subsection{Integration of a Weather Map System for Risk Monitoring and Early Warning}

% To provide both users and rescue teams with a comprehensive view of disaster developments, another important future enhancement is the integration of a visual weather map system. This system is intended to support real-time monitoring of critical environmental indicators, such as rainfall intensity and temperature distribution. To ensure high performance and scalability, the data collection and distribution pipeline will be designed using a multi-layer optimization approach.

% On the backend side, instead of allowing the main NestJS application to directly query third-party weather platforms (e.g., OpenWeather) whenever a request is received—a design that may introduce processing bottlenecks under high traffic conditions—a dedicated background worker will be deployed. This worker will periodically retrieve weather data from OpenWeather and synchronize the collected information into Redis, a high-performance in-memory data store. Consequently, when a user requests weather information, the NestJS server can simply fetch the required data from the Redis cache and return it immediately, significantly reducing network latency and minimizing the computational burden placed on the core application.

% On the mobile application side, the weather data processing workflow will be implemented in accordance with the Model–View–ViewModel (MVVM) architectural pattern to ensure efficient rendering and maintainable code organization. Specifically, when a user interacts with the interface to display rainfall or temperature maps, the ViewModel layer will initiate requests for weather map data and interpolated measurement points through the Repository layer. Acting as the system's single source of truth, the Repository will evaluate the validity and freshness of the available data to determine whether the information can be retrieved directly from the local cache or whether a network request must be sent to the NestJS backend. Once the raw weather dataset has been obtained, the Repository forwards the information to the ViewModel, where data transformation and presentation logic are applied. Finally, the View layer receives the processed data and renders the corresponding heatmap overlays together with relevant weather indicators on the device screen. This architecture not only provides users with a smooth and responsive interactive experience but also delivers timely environmental information that can support disaster awareness, risk assessment, and emergency preparedness activities.

% \subsection{Integration of a VietQR-Based Donation System for Automated Transaction Tracking}

% To enhance the automated financial tracking and reconciliation capabilities of the charity management subsystem, the future roadmap includes the complete integration of a VietQR-based donation system. At present, the core codebase for this feature has been developed and is currently in the testing phase. However, due to time constraints during the thesis implementation, it has not yet been deployed into production as several validation bugs and transaction synchronization edge cases are still being resolved.

% To support the transaction monitoring workflow, the PostgreSQL database schema will be updated, specifically requiring the addition of four new fields to the \texttt{User} table: \texttt{username\_from\_vietqr}, \texttt{password\_from\_vietqr}, \texttt{\seqsplit{username\_to\_vietqr}}, and \texttt{password\_to\_vietqr}. To ensure strict security compliance and protect sensitive API integration credentials, the password fields will be encrypted inside the database using two-way Advanced Encryption Standard (AES) encryption. Completing this integration will allow the system to automatically match bank transfer callbacks with charity campaign donation records, enhancing financial transparency without manual verification.

% \subsection{Integration of Delay-Tolerant Networking (DTN) for Disconnected Emergency Communication}

% To ensure operational continuity in catastrophic disaster scenarios where physical cellular towers and local Internet infrastructures are completely destroyed, a key future development is the integration of Delay-Tolerant Networking (DTN) technologies. Inspired by the Cellular Data Offloading System (Cedos) architecture, this enhancement aims to support opportunistic, store-carry-and-forward communication when end-to-end network paths are unavailable.

% On the mobile client side, the FloodHelper application will implement a \seqsplit{Delay/Disruption-Tolerant} Transport Protocol (D2TP) abstraction layer running in user-space over UDP. When a user loses internet connectivity, critical SOS signals, survival metadata, and tracking logs will be serialized and cached locally. The application will leverage opportunistic peer-to-peer connections (such as Wi-Fi Direct or Bluetooth Low Energy) to transmit these cached packets to moving nodes—such as rescue vehicles, drones, or other active users acting as "data mules"—that pass through the disconnected zone.

% To manage multiple concurrent data transfers efficiently on resource-constrained devices, a multi-flow scheduler (resembling the D2TS Earliest Deadline First model) will be developed to prioritize safety-critical SOS alerts over periodic location logs. Additionally, to prevent battery drainage during prolonged offline states, an active energy management mechanism based on the D-Lock waking protocol will be implemented. This will allow the device to perform brief, periodic scans for neighboring peer nodes while keeping the CPU and radio interfaces in a deep-sleep state for the majority of the time. Once a data mule encounters an active gateway or proxy (similar to D2Prox) connected to the Internet, the stored messages will be automatically synchronized with the central NestJS backend, thereby establishing a resilient, decentralized communications pipeline for severe flood conditions.

\end{document}